\documentclass[14pt,a4paper,oneside]{extbook}
\usepackage[utf8]{vietnam}
\usepackage{amsmath,amssymb,mathrsfs,fancyhdr,enumerate,multirow,makecell}
\usepackage{tkz-euclide,tikz-3dplot,tikz,tkz-tab}
\usetikzlibrary{arrows.meta,calc,intersections,angles,quotes,backgrounds,shapes.geometric,patterns,through,shadows,positioning}
\usepackage{yhmath}
\usepackage[top=0.7cm, bottom=0.7cm, left=1.0cm, right=1.0cm, headheight=16pt, headsep=3mm, footskip=5mm]{geometry}
\usepackage[hidelinks,unicode]{hyperref}
\usepackage{qrcode}
\usepackage[many]{tcolorbox}
\usepackage{fontawesome}

\pagestyle{fancy}
\fancyhf{}
\fancyhead[L]{\footnotesize\bfseries\color{teal!85!black}\faBook~SỔ TAY KIẾN THỨC TOÁN 4}
\fancyhead[R]{\footnotesize\bfseries\color{blue!85!black}Lớp Toán Thầy Lĩnh -- loptoanthaylinh.vn}
\fancyfoot[C]{\footnotesize\bfseries\color{gray!80!black}-- Sổ tay kiến thức: Bài 31. Hình bình hành, hình thoi --}
\renewcommand{\headrulewidth}{0.4pt}

\begin{document}

% --- BANNER TIÊU ĐỀ BÀI HỌC ---
\noindent
\begin{tcolorbox}[
  enhanced,
  colback=teal!8!white,
  colframe=teal!80!black,
  arc=3mm, boxrule=1pt,
  left=3mm, right=3mm, top=2mm, bottom=2mm,
  halign=center
]
  {\Large\bfseries\color{teal!90!black}\faBookmark\ BÀI 31: HÌNH BÌNH HÀNH, HÌNH THOI}\\[1.5pt]
  {\footnotesize\itshape\color{black!75} Đặc điểm cạnh đối song song, dấu hiệu phân biệt và tính chất hai đường chéo hình thoi trên 1 trang A4}
\end{tcolorbox}

\vspace{2mm}

% --- KHỐI 1 & 2: LÝ THUYẾT & KĨ NĂNG ---
\noindent
\begin{minipage}[t]{0.49\textwidth}
  \begin{tcolorbox}[
    enhanced,
    title={\footnotesize\bfseries\color{white}\faCube\ 1. Hình bình hành},
    colback=blue!2!white,
    colframe=blue!75!black,
    arc=2.5mm, boxrule=0.8pt, top=1.8mm, bottom=1.8mm, left=2mm, right=2mm
  ]
    \footnotesize
    $\bullet$ \textbf{Đặc điểm}: Hình bình hành có \textbf{hai cặp cạnh đối diện song song và bằng nhau}:
    \begin{center}
      $AB \parallel CD$ và $AB = CD$\\
      $AD \parallel BC$ và $AD = BC$
    \end{center}
    $\bullet$ Hai góc đối diện bằng nhau.\\
    $\bullet$ \textbf{Hình ảnh}: Lan can cầu thang, khung lưới bóng đá, hoa văn trang trí.
  \end{tcolorbox}
\end{minipage}%
\hfill
\begin{minipage}[t]{0.49\textwidth}
  \begin{tcolorbox}[
    enhanced,
    title={\footnotesize\bfseries\color{white}\faDiamond\ 2. Hình thoi},
    colback=teal!2!white,
    colframe=teal!75!black,
    arc=2.5mm, boxrule=0.8pt, top=1.8mm, bottom=1.8mm, left=2mm, right=2mm
  ]
    \footnotesize
    $\bullet$ \textbf{Đặc điểm}: Hình thoi có hai cặp cạnh đối diện song song và \textbf{bốn cạnh bằng nhau}:
    \begin{center}
      $AB = BC = CD = DA$
    \end{center}
    $\bullet$ \textbf{Tính chất đặc biệt 2 đường chéo}:\\
    --- \textbf{Vuông góc} với nhau ($AC \perp BD$).\\
    --- \textbf{Cắt nhau tại trung điểm} của mỗi đường.\\
    $\bullet$ \textbf{Hình thoi là hình bình hành đặc biệt}.
  \end{tcolorbox}
\end{minipage}

\vspace{2mm}

% --- KHỐI 3: SƠ ĐỒ TƯ DUY TIKZ ---
\noindent
\begin{tcolorbox}[
  enhanced,
  title={\footnotesize\bfseries\color{white}\faSitemap\ SƠ ĐỒ TƯ DUY: BẢNG SO SÁNH TRỰC QUAN HÌNH BÌNH HÀNH VÀ HÌNH THOI},
  colback=yellow!2!white,
  colframe=orange!85!black,
  arc=2.5mm, boxrule=0.8pt, top=1.5mm, bottom=1.5mm, left=1mm, right=1mm,
  halign=center
]
  \begin{tikzpicture}[
    scale=0.92,
    box/.style={rectangle, rounded corners=2mm, draw=teal!80!black, fill=teal!10, thick, inner sep=3pt, font=\scriptsize\bfseries, align=center}
  ]
    % Hình bình hành ABCD
    \coordinate (A) at (0.8, 0.9);
    \coordinate (B) at (3.5, 0.9);
    \coordinate (C) at (2.7, -0.6);
    \coordinate (D) at (0.0, -0.6);

    \draw[thick, blue!85!black, fill=blue!8] (A) -- (B) -- (C) -- (D) -- cycle;
    \node[above left, font=\tiny\bfseries] at (A) {$A$};
    \node[above right, font=\tiny\bfseries] at (B) {$B$};
    \node[below right, font=\tiny\bfseries] at (C) {$C$};
    \node[below left, font=\tiny\bfseries] at (D) {$D$};

    \node[box, draw=blue!75!black, fill=blue!8, text width=5.2cm] at (1.8, -1.35) {\textbf{HÌNH BÌNH HÀNH}\\2 cặp cạnh đối song song\\$AB = CD$, $AD = BC$};

    % Dấu phân cách
    \draw[thick, dashed, gray!60] (4.8, 1.2) -- (4.8, -1.8);

    % Hình thoi MNPQ
    \coordinate (M) at (7.8, 1.0);
    \coordinate (N) at (9.6, 0.15);
    \coordinate (P) at (7.8, -0.7);
    \coordinate (Q) at (6.0, 0.15);

    \draw[thick, teal!85!black, fill=teal!8] (M) -- (N) -- (P) -- (Q) -- cycle;
    \draw[thick, dashed, red!80!black] (M) -- (P);
    \draw[thick, dashed, red!80!black] (N) -- (Q);
    \node[font=\tiny\bfseries] at (7.8, 0.15) {$O$};
    \draw[thick, red!80!black] (7.8, 0.35) -- (8.0, 0.35) -- (8.0, 0.15);

    \node[above, font=\tiny\bfseries] at (M) {$M$};
    \node[right, font=\tiny\bfseries] at (N) {$N$};
    \node[below, font=\tiny\bfseries] at (P) {$P$};
    \node[left, font=\tiny\bfseries] at (Q) {$Q$};

    \node[box, draw=teal!80!black, fill=teal!8, text width=5.5cm] at (7.8, -1.35) {\textbf{HÌNH THOI}\\4 cạnh bằng nhau: $MN=NP=PQ=QM$\\2 đường chéo: \textbf{VUÔNG GÓC} tại $O$};

    % Thẻ tổng kết
    \node[box, draw=purple!80!black, fill=purple!10, text width=4.5cm] at (13.5, 0.0) {\textbf{\color{purple!85!black}MỐI QUAN HỆ CỐT LÕI}\\[1mm]Hình thoi chính là \textbf{hình bình hành đặc biệt}\\có 4 cạnh bằng nhau và 2 đường chéo vuông góc};
  \end{tikzpicture}
\end{tcolorbox}

\vspace{2mm}

% --- KHỐI 4: VÍ DỤ MINH HỌA & BẪY TRẮC NGHIỆM ---
\noindent
\begin{minipage}[t]{0.53\textwidth}
  \begin{tcolorbox}[
    enhanced,
    title={\footnotesize\bfseries\color{white}\faCheckCircle\ Ví dụ thực hành chuẩn mực},
    colback=green!2!white,
    colframe=green!60!black,
    arc=2.5mm, boxrule=0.8pt, top=1.8mm, bottom=1.8mm, left=2.2mm, right=2.2mm
  ]
    \footnotesize
    $\bullet$ Cho hình thoi $ABCD$ có cạnh $AB = 6\text{ cm}$. Tính chu vi hình thoi:\\
    Vì hình thoi có 4 cạnh bằng nhau nên:\\
    $\text{Chu vi} = 6 \times 4 = \mathbf{24\text{ cm}}$.\\[1mm]
    $\bullet$ Hai đường chéo $AC$ và $BD$ cắt nhau tại $O$. Biết $AC = 8\text{ cm}$, $BD = 6\text{ cm}$:\\
    $O$ là trung điểm nên $OA = OC = 4\text{ cm}$; $OB = OD = 3\text{ cm}$ và $\widehat{AOB} = 90^\circ$.
  \end{tcolorbox}
\end{minipage}%
\hfill
\begin{minipage}[t]{0.45\textwidth}
  \begin{tcolorbox}[
    enhanced,
    title={\footnotesize\bfseries\color{white}\faExclamationTriangle\ Mẹo tránh bẫy trắc nghiệm},
    colback=red!2!white,
    colframe=red!70!black,
    arc=2.5mm, boxrule=0.8pt, top=1.8mm, bottom=1.8mm, left=2.2mm, right=2.2mm
  ]
    \footnotesize
    $\bullet$ \textbf{Bẫy hình thoi và hình vuông:}\\
    Hình thoi \textbf{chưa chắc} là hình vuông (vì 4 góc chưa chắc vuông). Chỉ khi hình thoi có 4 góc vuông thì mới là hình vuông!\\
    $\bullet$ \textbf{Bẫy đường chéo hình bình hành}: Đường chéo hình bình hành thường \textbf{không vuông góc} với nhau!
  \end{tcolorbox}
\end{minipage}

\end{document}
